% Requirement
% Motivation
% What is the general research topic?
% What is the significance of the research topic? Why is this topic worth working on?
% Literature Review
% Provide a broad view of the relevant research landscape
% Provide concrete descriptions on a small handful of the most highly related works.
% Identify major trends/patterns across prior work.
% Identify shortcomings with prior work and/or other opportunities that motivate this work
% Initial Proposed Approaches
% Provide high-level descriptions for at least 2 potential approaches to innovate on the research topic
% see "Suggestions for Brainstorming Project Ideas" below for concrete advice
% be sure to consider compute requirements, and provide an estimate in your report
% Initial Proposed Data
% Identify data sources that will be useful for training (if applicable) and evaluation.
% Initial Proposed Evaluation
% Which concepts or theoretical constructs are most relevant for evaluating your method, and why?
% If the constructs can’t be directly measured, what metrics will serve as acceptable proxies?
% If your system consists of multiple components, what constructs/metrics will you use to evaluate the sub-components?
% What qualitative analyses do you plan to conduct?



% \begin{abstract}
% Abstract here
% \end{abstract}

\section{Motivation} % Introduction

% 1. Structure decoding importance
% The difference in diffusion and AR structure decoding
% why it worth research for structure decoding in diffusion
% Current downside
% constrained decoding + LAVE
% their downside O(N) pick candidates
% solve: CSR matrix + Herding algorithm

% 2. Deterministic > probabilistic
% sampling wall causing information collapse

Diffusion large language models (dLLMs) have emerged as a promising alternative to autoregressive (AR) models, offering parallel token generation and bidirectional context modeling. A critical challenge in deploying dLLMs for structured generation tasks, as shown in figure \ref{fig:constraint-pass-fail}, such as JSON schema compliance, code synthesis, and chemical expression generation, is ensuring that outputs reliably conform to formal grammar constraints \citep{cd4dllm}.

\begin{figure*}[t]
\centering
\resizebox{\textwidth}{!}{%
    \begin{tikzpicture}[
        font=\small,
        tok/.style={
            draw=black,
            line width=0.5pt,
            minimum width=0.78cm,
            minimum height=0.52cm,
            align=center,
            inner sep=1pt
        },
        good/.style={
            tok,
            fill=blue!10
        },
        bad/.style={
            tok,
            fill=red!10
        },
        title/.style={
            font=\small\bfseries
        },
        noteok/.style={
            font=\small\bfseries,
            text=green!50!black
        },
        notebad/.style={
            font=\small\bfseries,
            text=red!70!black
        },
        smallnote/.style={
            font=\scriptsize
        },
        arr/.style={
            -{Stealth[length=2.2mm,width=1.6mm]},
            line width=0.9pt,
            draw=black!70
        }
    ]
    
    % =========================
    % Left panel
    % =========================
    \node[title] at (3.2,1.8) {(a) Without constrained decoding};
    
    \node[tok] (L1) at (0,0) {\{};
    \node[tok, right=0cm of L1] (L2) {"name"};
    \node[tok, right=0cm of L2] (L3) {:};
    \node[tok, right=0cm of L3] (L4) {"Bob"};
    \node[bad, right=0cm of L4] (L5) {"age"};
    \node[tok, right=0cm of L5] (L6) {:};
    \node[tok, right=0cm of L6] (L7) {30};
    \node[bad, right=0cm of L7] (L8) {\}};
    
    \node[smallnote, text=red!70!black] at ($(L5.north)+(0,0.45)$) {missing comma};
    
    \node[notebad] at ($(L1.south)!0.5!(L8.south)+(0,-0.5)$) {Grammar checker: fail};
    
    % =========================
    % Right panel
    % =========================
    \begin{scope}[xshift=9.8cm]
    
    \node[title] at (3.45,1.8) {(b) With constrained decoding};
    
    \node[tok] (R1) at (0,0) {\{};
    \node[tok, right=0cm of R1] (R2) {"name"};
    \node[tok, right=0cm of R2] (R3) {:};
    \node[good, right=0cm of R3] (R4) {"Bob"};
    \node[good, right=0cm of R4] (R5) {,};
    \node[tok, right=0cm of R5] (R6) {"age"};
    \node[tok, right=0cm of R6] (R7) {:};
    \node[good, right=0cm of R7] (R8) {30};
    \node[tok, right=0cm of R8] (R9) {\}};
    
    \draw[arr] ($(R5.north)+(0,0.14)$)
        .. controls +(0,0.42) and +(0.18,0.30) ..
        ($(R5.north)+(0.02,0.02)$);
    
    \node[smallnote, text=black!70] at ($(R5.north)+(0.15,0.72)$) {valid token inserted};
    
    \node[noteok] at ($(R1.south)!0.5!(R9.south)+(0,-0.5)$) {Grammar checker: pass};
    
    \end{scope}
    
    \end{tikzpicture}%
}
\caption{
Example of unconstrained versus grammar-constrained decoding.
Without constraints, the model may emit an invalid continuation, causing the grammar check to fail.
With constrained decoding, invalid continuations are pruned during generation, yielding a grammar-valid output.
}
\label{fig:constraint-pass-fail}
\end{figure*}


The state-of-the-art approach, LAVE \citep{zhang2026lookaheadthenverifyreliableconstraineddecoding}, addresses this by leveraging dLLMs' ability to predict token distributions for all positions in parallel during each forward pass. When the model proposes a token, LAVE performs lookahead by sampling $N$ complete prefixes for the remaining masked positions and accepts the token only if at least one sampled prefix passes a grammar parser. While effective, this approach has two fundamental limitations.

First, the verification process in LAVE is inherently stochastic. Because it relies on random sampling to fill masked positions, there is a persistent, non-zero probability that a valid token will be erroneously rejected. This occurs when none of the $N$ generated samples yield a parseable prefix, triggering unnecessary retries and degraded generation quality.

Second, LAVE evaluates token decisions independently at each position, mirroring the left-to-right bottlenecks of AR models. When multiple positions are unmasked during a single diffusion step, LAVE verifies each token in isolation, ignoring the joint dependencies between simultaneously unmasked slots. Thus, while individual tokens may pass local verification, their combination can create a dead-end for subsequent tokens, rendering them illegal. This greedy, position-by-position strategy is fundamentally misaligned with the parallel nature of diffusion generation.

We propose Grammar-Guided Beam Search (GGBS), a deterministic constrained decoding framework that integrates structural constraints directly into the beam expansion process. By replacing post-hoc random sampling with incremental parser state tracking, GGBS eliminates the stochastic rejection of valid tokens; instead, it performs a deterministic top-$K$ selection from a set of grammar-valid candidates. Crucially, GGBS shifts the decoding paradigm from independent token validation to a joint optimization over all masked positions within a diffusion step. By maintaining parser states across both fixed and masked segments, GGBS identifies the highest-probability token assignment that satisfies global grammar constraints across the entire unmasked span. This approach aligns the decoding process with the parallel nature of diffusion, ensuring structural coherence without the efficiency trade-offs of greedy, position-by-position verification. The code is available at github \url{https://github.com/buffett0323/anlp_final}.


% While discrete diffusion language models (dLLMs) offer a parallelized alternative to autoregressive (AR) generation, our analysis identifies a Trilemma of Inefficiency that prevents their deployment in production-scale structured decoding.

% First, the system is hindered by \textbf{search width}, where enforcing constraints requires an exhaustive $O(V)$ scan of the entire vocabulary to identify valid transitions at every step. Existing constrained diffusion frameworks, such as DINGO \citep{suresh2025dingoconstrainedinferencediffusion}, require an exhaustive scan over the full vocabulary $V$ to identify valid transitions at every denoising step. This linear dependence on vocabulary size introduces substantial memory and I/O overhead on modern accelerators, undermining the throughput advantage of non-autoregressive generation.

% Second, it also suffers from a lack of \textbf{information integrity} due to the ``sampling wall,'' a phenomenon where rich probability distributions are discarded in favor of sparse one-hot token commitments \citep{jo2026loopholingdiscretediffusiondeterministic}. Standard dLLMs rely on stochastic ancestral sampling, in which rich token distributions are collapsed into one-hot commitments at each step. In structured generation, a high-probability token may be temporarily disallowed by a grammar constraint (e.g., a required JSON bracket). Once suppressed, that latent preference can be lost, causing the model to forget its original semantic intent and increasing the risk of logical dead-ends.

% Third, the \textbf{inference depth} remains a bottleneck, as standard unmasking schedules require a high number of function evaluations (NFE)\footnote{NFE is the number of times the model must perform a forward denoising step during inference to produce an output.} to maintain syntactic accuracy, often nullifying the theoretical latency benefits inherent to the diffusion paradigm. High-quality structured generation often requires a large NFE. Although speculative decoding has become an effective acceleration technique for AR models \citep{leviathan2023fastinferencetransformersspeculative}, dLLMs still lack a mechanism for verifying multiple future structural positions within a single forward pass while preserving strict syntactic correctness.

% To address these barriers, we propose the \textbf{Sparse Deterministic Speculative Decoding (SD2)} framework. SD2 combines three key ideas: (i) flattening formal automata into CSR matrices for effectively $O(1)$ vectorized token retrieval; (ii) replacing stochastic sampling with a deterministic herding algorithm that preserves token momentum across denoising steps; and (iii) introducing grammar-legal speculative trees that exploit the bidirectional context of dLLMs to accelerate convergence. By combining hardware-native indexing, deterministic intent preservation, and structure-aware speculation, SD2 provides a scalable pathway for deploying strictly constrained diffusion engines in latency-sensitive environments.


% -----

% While discrete diffusion language models (dLLMs) offer a parallelized alternative to autoregressive (AR) generation, our analysis shows that the current state of the art still faces three critical barriers to production-scale structured decoding. First, the system is hindered by \textbf{Search Width}, where enforcing constraints requires an exhaustive $O(N)$ scan of the entire vocabulary to identify valid transitions at every step. Second, it suffers from a lack of \textbf{Information Integrity} due to the ``sampling wall,'' a phenomenon where rich probability distributions are discarded in favor of sparse one-hot token commitments. Finally, the \textbf{Inference Depth} remains a bottleneck, as standard unmasking schedules require a high number of function evaluations (NFE) to maintain syntactic accuracy, often nullifying the theoretical latency benefits inherent to the diffusion paradigm.

% To address these challenges, existing works have proposed mathematically rigorous but computationally expensive solutions. DINGO \citep{suresh2025dingoconstrainedinferencediffusion} successfully introduced the first provably correct dynamic programming (DP) framework for dLLMs to ensure 100\% syntactic correctness. However, its transition cost calculation iterates through the full vocabulary $V$, resulting in a complexity of $O(d \cdot |Q| \cdot |V|)$ per block. This linear dependence on vocabulary size creates a significant I/O overhead on modern accelerators. Furthermore, the reliance on stochastic ancestral sampling leads to a ``sampling wall'' \citep{jo2026loopholingdiscretediffusiondeterministic}, where the soft information about alternative token candidates collapses into a single category. In structured tasks, where high-probability tokens are often suppressed by temporary grammar constraints (e.g., a required JSON bracket), this collapse prevents the model from remembering its original intent, leading to oscillations or logical dead-ends. While recent paradigms like LAVE \citep{zhang2026lookaheadthenverifyreliableconstraineddecoding} and Plan-Verify-Fill (PVF) \citep{li2026planverifyfillstructured} attempt to improve reliability through lookahead or structural anchors, they primarily serve as statistical filters and do not resolve the underlying $O(N)$ retrieval complexity or the serial nature of step-by-step unmasking.

% We propose the Sparse Deterministic Speculative Decoding (SD2) framework to systematically overcome these three dimensions of inefficiency. To optimize the search width, we incorporate the STATIC framework \citep{su2026vectorizingtrieefficientconstrained} to flatten formal automata into hardware-friendly Compressed Sparse Row (CSR) matrices, recasting token retrieval as vectorized $O(1)$ memory access through coalesced reads. To ensure information integrity, we replace stochastic sampling with a deterministic Herding algorithm \citep{herding}, maintaining a continuous weight vector that preserves the momentum of suppressed tokens across denoising steps, allowing blocked intents to resurface once the grammar state permits \citep{suzuki2026deterministicdiscretedenoising}. Finally, we address inference depth by constructing grammar-legal speculative trees, leveraging the inherent bidirectional context of dLLMs to verify multiple future structural positions in a single forward pass. By synergizing sub-linear sparse indexing with deterministic intent preservation, SD2 provides a scalable pathway for deploying strictly constrained diffusion engines in real-world, latency-sensitive environments such as automated programming and symbolic reasoning. The code is available at github \url{https://github.com/buffett0323/SD2}.



\section{Literature Review}

\subsection{Discrete Diffusion and Parallel Decoding}
Discrete diffusion models (dLLMs) such as LLaDA \citep{llada} and Dream \citep{dream} represent a paradigm shift from traditional AR models by enabling non-autoregressive, any-order text generation through iterative refinement. However, this flexibility introduces the incomplete prefix problem: unlike AR models, where a stable prefix always exists for grammar state tracking, dLLMs must often predict tokens surrounded by unmasked or partially denoised noise. While recent acceleration techniques like DFlash \citep{chen2026dflashblockdiffusionflash} and Spiffy \citep{agrawal2026spiffymultiplyingdiffusionllm} reduce the sampling steps via block-diffusion, they remain ``unconstrained,'' failing to guarantee structural or syntactic correctness in the generated output.

\subsection{Constrained Decoding: From AR to dLLMs}
Traditional constrained decoding frameworks like \textit{Outlines} \citep{outlines} and \textit{XGrammar} \citep{dong2024xgrammar} rely on the linear progression of AR decoding to maintain finite state machine (FSM) or parser states. These methods cannot be directly applied to dLLMs due to the lack of a monotonic causal chain. 
Transitioning these constraints to the diffusion paradigm, \citet{cd4dllm} proposed the first Context-Free Grammer (CFG-constrained) dLLM decoder, while DINGO \citep{suresh2025dingoconstrainedinferencediffusion} utilized dynamic programming (DP) to ensure syntactic correctness, albeit at a high computational cost. LAVE \citep{zhang2026lookaheadthenverifyreliableconstraineddecoding} introduced a ``lookahead-then-verify'' mechanism as a baseline for reliability; however, LAVE's reliance on stochastic sampling for verification possibly lead to the rejection of valid tokens. Our work draws inspiration from the STATIC framework \citep{su2026vectorizingtrieefficientconstrained}, which demonstrated that sparse transition structures can be flattened into Compressed Sparse Row (CSR) matrices to enable $O(K)$ coalesced lookups, where $K$ is the branch factor, rather than scanning the full vocabulary of size $|V|$. We adapt this principle to represent DFA token-level transitions for grammar-constrained decoding, reducing the per-step transition cost in our bidirectional Viterbi DP from $O(|V|)$ to $O(K)$.


% \subsection{Deterministic Dynamics in Discrete Spaces}
% A significant challenge in discrete diffusion is the "sampling wall"—the information loss occurring when categorical distributions are discretized via random sampling. Research into Deterministic Discrete Denoising \citep{suzuki2026deterministicdiscretedenoising} and the Herding algorithm \citep{herding} suggests that replacing stochastic sampling with deterministic state transitions can preserve the "working memory" of suppressed tokens. By adopting deterministic dynamics, our proposed GGBS eliminates the "luck-of-the-draw" failure modes found in LAVE, ensuring that the highest-probability grammar-valid paths are consistently selected rather than sampled.

\subsection{Beam Search and Joint Multi-Token Decoding}
While classic beam search is a staple of AR decoding, its application in dLLMs requires a fundamental shift toward joint optimization. In diffusion steps where multiple positions are unmasked simultaneously, independent token-wise verification fails to account for inter-token dependencies, often leading to globally invalid sequences. \citet{leviathan2023fastinferencetransformersspeculative} has demonstrated the power of parallel verification in AR contexts; however, existing dLLM methods often revert to greedy, position-by-position checks. GGBS addresses this by treating the unmasked span as a joint decoding problem, utilizing bidirectional suffix feasibility pruning to ensure that the chosen tokens satisfy grammar constraints globally rather than just locally.


% \section{Literature Review}

% \subsection{Discrete Diffusion and Parallel Decoding}
% Discrete diffusion models like LLaDA-8B \citep{llada} and Dream-7B \citep{dream} generate text through iterative refinement, breaking the strict causal constraints of AR models. While they enable any-order generation, their practical utility is often limited by high NFE requirements. Recent efforts like DFlash \citep{chen2026dflashblockdiffusionflash} and Spiffy \citep{agrawal2026spiffymultiplyingdiffusionllm} attempt to reduce this overhead via block-diffusion or self-speculative mechanisms, yet these acceleration techniques have not been natively integrated with formal structural constraints.

% \subsection{Formal Constraints in dLLMs}
% The transition from statistical generation to symbolic accuracy has led to frameworks like DINGO \citep{suresh2025dingoconstrainedinferencediffusion}, which employs dynamic programming (DP) to ensure 100\% syntactic correctness in diffusion. However, the computational cost of DP in dLLMs remains high. While LAVE \citep{zhang2026lookaheadthenverifyreliableconstraineddecoding} and PVF \citep{li2026planverifyfillstructured} improve reliability through lookahead or planning strategies, they function primarily as high-level filters. Our work builds upon the STATIC framework \citep{su2026vectorizingtrieefficientconstrained}, which demonstrated that flattening trie-based constraints into Compressed Sparse Row (CSR) matrices can transform irregular memory access into hardware-friendly vectorized operations, offering a path to bypass the $O(|V|)$ retrieval bottleneck.

% \subsection{Deterministic Dynamics and Intent Preservation}
% A core challenge in discrete state spaces is the ``sampling wall'', the information loss occurring when categorical distributions are discretized. Research into Deterministic Discrete Denoising \citep{suzuki2026deterministicdiscretedenoising} suggests that transforming stochastic sampling into deterministic state transitions can maintain a ``working memory'' of suppressed tokens. This approach, rooted in the Herding algorithm, allows models to achieve faster asymptotic convergence. SD2 integrates these deterministic dynamics specifically to solve the intent-grammar mismatch often found in structured decoding.

% \subsection{Speculative Tree Verification}
% Speculative decoding has successfully accelerated AR models by using draft models to propose multiple tokens for parallel verification \citep{leviathan2023fastinferencetransformersspeculative}. In the diffusion paradigm, SSD leverages the model’s inherent parallel prediction capability. However, existing speculative methods for diffusion often revert to expensive pointer-chasing structures for constraint enforcement. SD2 extends this by constructing speculative trees that are inherently grammar-legal, utilizing the bidirectional nature of dLLMs to verify structural anchors across multiple positions simultaneously.


% \subsection{Discrete Diffusion Language Models}
% Discrete diffusion models have emerged as a significant alternative to AR architectures, enabling text generation through iterative refinement in discrete state spaces. Unlike AR models that are bound by strict causal factorization, models such as LLaDA-8B \citep{llada} and Dream-7B \citep{dream} utilize bidirectional attention to decode token blocks in parallel, offering superior flexibility for any-order generation and infilling. Despite their parallel potential, these models often require a high NFE to match the quality of AR counterparts, motivating the need for more efficient unmasking and remasking strategies.

% \subsection{Constrained Inference in Non-Autoregressive Paradigms}
% Enforcing structural correctness in probabilistic generators has traditionally relied on finite-state machines (FSM) or pushdown automata (PDA) to mask logits during sequential decoding, as seen in libraries like Outlines \citep{outlines} and XGrammar \citep{dong2024xgrammar}. DINGO \citep{suresh2025dingoconstrainedinferencediffusion} represents the first provably correct constrained decoding framework for diffusion LLMs, employing DP to find the highest-probability path that satisfies regular expression constraints. However, DINGO’s search process scales with the full vocabulary size $N$, incurring significant runtime overhead. More recently, LAVE \citep{li2026planverifyfillstructured} introduced a "lookahead-then-verify" mechanism to improve reliability for context-free grammars (CFG), yet it primarily serves as a statistical filter without addressing the underlying computational complexity of token retrieval.


% \subsection{Deterministic Dynamics and the Sampling Wall}
% A fundamental limitation of discrete diffusion is the ``sampling wall,'' where rich categorical distributions collapse into sparse one-hot vectors, leading to information loss across denoising steps. \citet{suzuki2026deterministicdiscretedenoising} proposed Deterministic Discrete Denoising based on the herding algorithm, which transforms stochastic sampling into deterministic state transitions governed by continuous weight vectors. This mechanism effectively maintains a working memory of suppressed token momentum, allowing for faster asymptotic convergence ($O(T^{-1})$) and higher sample quality. Recent paradigms like PVF \citep{li2026planverifyfillstructured} also advocate for active structural planning over reactive harvesting, but they have yet to integrate $O(1)$ sparse indexing for the sub-linear optimization of formal grammar transitions.


% \subsection{Acceleration via Speculative and Parallel Decoding}
% To alleviate the inference bottleneck, speculative decoding utilizes a lightweight draft model to propose tokens that a larger target model verifies in parallel. In the diffusion context, Self-Speculative Decoding (SSD) \citep{gao2025selfspeculativedecodingdiffusion} and Spiffy \citep{agrawal2026spiffymultiplyingdiffusionllm} leverage the model’s inherent parallel prediction capability to act as its own drafter, accelerating inference without auxiliary modules. Hybrid approaches like DFlash \citep{chen2026dflashblockdiffusionflash} further improve throughput by using lightweight block-diffusion drafters conditioned on target hidden states. While these methods reduce NFE, they lack native mechanisms to enforce hardware-efficient structural constraints during the drafting phase, often reverting to expensive pointer-chasing trie structures.

% \subsection{Hardware-Native Optimization with STATIC}
% Modern hardware accelerators are severely penalized by the irregular memory access patterns of traditional prefix trees (tries). The STATIC \citep{su2026vectorizingtrieefficientconstrained} framework addresses this by flattening trie-specified constraints into CSR matrices, recasting graph traversals as vectorized sparse matrix operations. By achieving $O(1)$ I/O complexity with respect to the constraint set size, STATIC enables massive efficiency gains (up to 948x) on GPUs and TPUs. Integrating these hardware-native indexing principles with herding-based deterministic unmasking provides a novel pathway to break the $O(N)$ search barrier in structured diffusion decoding.



\section{Initial Approach: Grammar-Guided Beam Search (GGBS)}

We adopt \emph{grammar-guided beam search} as our initial constrained-decoding strategy for masked diffusion models. The core idea is to treat each bidirectional gap of mask positions as a short decoding horizon: instead of committing to a single argmax token per step, we expand a bounded beam of hypotheses whose next-token choices are filtered by a token-level grammar checker. This yields joint decisions over multiple mask columns while keeping the same forward passes as standard decoding.

\paragraph{Mechanics.}
At a gap frontier, we maintain beams indexed by partial token assignments over the masked span. For each beam we track the parser state implied by the committed prefix (prompt plus non-mask tokens plus previously chosen gap tokens). At each mask column we score candidates from model logits (typically top-$K$ by probability) and retain only those that the parser can consume in one step; we then rank beams by cumulative log-probability and prune to width $B$. When a fixed (non-mask) segment separates two mask columns, we advance the parser along those tokens without beam branching. Optionally, we enforce suffix feasibility on the last mask before a fixed right context (e.g., closing JSON) so that beams that cannot extend to the remainder are discarded early. The detail is shown in algorithm \ref{alg:ggbs}.


\begin{algorithm*}[t]
\caption{Grammar-Guided Beam Search (GGBS) over a bidirectional mask gap}
\label{alg:ggbs}
\begin{algorithmic}[1]
\Require Prompt tokens, diffusion state with mask gap positions
  $\mathcal{M}=\{m_1,\ldots,m_M\}$ in order; model logits on mask rows;
  grammar checker with incremental consume / rollback; beam width $B$;
  candidate pool size $K$ per mask column.
\Ensure One hypothesis over $\mathcal{M}$ (or failure).
\State Initialize beam $\mathcal{B} \gets \{(\text{prefix},\ \log p \gets 0)\}$.
\For{each mask position $m_i \in \mathcal{M}$ in left-to-right order}
  \State $\mathcal{B}' \gets \emptyset$
  \For{each beam $(s,\log p)\in \mathcal{B}$}
    \State Advance parser to the state \emph{before} token $x_{m_i}$ under hypothesis $s$.
    \State Take top-$K$ token ids by model probability at position $m_i$.
    \For{each candidate token $v$ in \emph{top-}$K$}
      \If{$v$ is grammar-legal at the current parser state}
        \State $s' \gets s$ with $x_{m_i}\gets v$; $\log p' \gets \log p + \log \pi(v\mid m_i)$
        \State Add $(s',\log p')$ to $\mathcal{B}'$ (apply optional suffix filter at last mask).
      \EndIf
    \EndFor
  \EndFor
  \State $\mathcal{B} \gets$ top-$B$ beams in $\mathcal{B}'$ by $\log p$.
  \If{$\mathcal{B}=\emptyset$} \Return failure \EndIf
\EndFor
\State \Return best beam in $\mathcal{B}$ by $\log p$.
\end{algorithmic}
\end{algorithm*}

\paragraph{Positioning relative to verify-and-resample.}
This design differs from ``generate freely, then verify'' pipelines: grammar is used during expansion to prune invalid continuations, not only after a full block is sampled. The trade-off is extra CPU work per validation step (parser probes and, in some settings, beam-local state updates) versus fewer wasted samples and more systematic search over the gap. We use this approach as the starting point for our experiments, aiming to combine deterministic decoding with bidirectional and joint decoding: deterministic decoding helps reduce false rejections, while bidirectional/joint decoding mitigates suffix incompatibility and addresses the limitations of per-token independent verification.



% To resolve the trilemma of search width, information integrity, and inference depth, we propose the \textbf{SD2} framework. The core of our approach lies in three interconnected modules:

% \begin{itemize}
%     \item \textbf{Vectorized Sparse Indexing:} We move away from pointer-chasing trie structures. By adopting the \texttt{STATIC} framework, we compile the target grammar (JSON, Python, or SQL) into a CSR matrix. This allows the dLLM to perform token masking as a single-pass bitwise operation across the vocabulary, reducing the transition complexity from $O(|V|)$ to $O(1)$ per valid state via coalesced memory reads.
    
%     \item \textbf{Deterministic Denoising via Herding:} We implement a discrete herding algorithm to bridge the sampling wall. Instead of collapsing the logits into a one-hot vector at each step $t$, we maintain a continuous error-accumulation vector $\mathbf{e}_t$. This vector tracks the difference between the model's ideal distribution and the constrained discrete output, ensuring that suppressed token intents (tokens that are highly probable but currently grammar-illegal) are preserved and prioritized as soon as the grammar state allows.
    
%     \item \textbf{Grammar-Legal Speculative Trees:} We introduce a bi-directional verification mechanism where the model predicts a tree of potential future tokens. Unlike AR speculation, we utilize the dLLM's ability to unmask multiple non-sequential indices. Each branch of the tree is pre-filtered by our CSR-based constraints, allowing the model to verify an entire structural block (e.g., an entire key-value pair in JSON) in a single forward pass.
% \end{itemize}


\section{Initial Proposed Data}

We ground the study in JSON-Bench first, using its test split of 272 instances, as the structure is easier, where json dataset won't contain infinite recursive loop like C++ code. Each instance pairs a natural-language user request with a JSON Schema (and reference output) and is turned into a standard chat prompt that instructs the model to answer only with JSON obeying that schema. The model under study is LLaDA-8B-Instruct with a fixed generation length (e.g., 256 tokens) and the same diffusion settings across methods (e.g., T = 128 steps, block length 32), so that differences in results come from the constraint mechanism, not from the backbone or decoding length. In later experiments, we plan to extend the study to CFG-intensive benchmarks such as CPP-Bench or Smile-Bench, where the grammar may require stack-based state tracking for nested structures such as parentheses, leading to more challenging, potentially unbounded parsing behavior.



\section{Initial Proposed Evaluation}
We evaluate both the syntactic validity and functional correctness of the generated JSON outputs. Syntactic@1 measures whether the top-1 output is valid JSON and satisfies the target schema, following the ETH / JSON-Bench evaluation pipeline. Functional@1, when applicable, measures task success against the benchmark’s reference tests. We also report wall-clock latency per instance, and, in instrumented runs, optionally break this down into model forward time and constraint-processing time, all under identical diffusion hyperparameters. We compare our method against strong constrained-diffusion baselines, including LAVE mask-and-verify approaches and our proposed Grammar-Guided Beam Search (GGBS) variant. This setup allows us to evaluate not only output validity, but also the runtime trade-offs between grammar-guided search over masked gaps and verify-after-sampling designs.



\section{Baseline Reproduction}

\paragraph{Baselines.}
We reproduce three baselines for our structured JSON generation task using LLaDA-8B-Instruct:

\begin{itemize}
    \item \textbf{NO-CD} (No Constrained Decoding): The plain LLaDA-8B-Instruct baseline with zero grammar enforcement. The diffusion model runs normally without any validity checks, serving as the unconstrained lower bound.
    \item \textbf{IG-CD} (Incremental Grammar-Constrained Decoding): The ETH baseline method \citep{cd4dllm} from the constrained-diffusion repository. It enforces grammar constraints during diffusion decoding by incrementally checking grammar validity at each unmasking step.
    \item \textbf{LAVE} \citep{zhang2026lookaheadthenverifyreliableconstraineddecoding}: The current state-of-the-art ``lookahead-then-verify’’ method that leverages dLLMs’ parallel prediction to sample $N$ complete prefixes and accept a token only if at least one sampled prefix passes a grammar parser.
\end{itemize}

LAVE is not the absolute state-of-the-art on all tasks, but it represents the strongest open-source constrained-decoding baseline for dLLMs and is the most directly comparable method to our proposed GGBS approach. IG-CD is included as the original ETH baseline to illustrate the incremental grammar-checking paradigm that GGBS also builds upon.

\paragraph{Reproduction Results on JSON-Bench.}
All experiments use JSON-Bench (272 instances), LLaDA-8B-Instruct, temperature $T=0.2$, generation length 256 tokens, seed 0, and $T=128$ diffusion steps on an A100 GPU. Table~\ref{tab:main-comparison} reports syntactic validity, functional correctness, and wall-clock latency alongside the originally reported figures.

\begin{table*}[h]
\centering
\caption{Main comparison on JSON-Bench (272 instances). Syntactic and functional scores are Syntactic@1 and Functional@1. Constraint \% denotes the fraction of tokens affected by grammar enforcement.}
\label{tab:main-comparison}
\begin{tabular}{lcccccc}
\hline
\textbf{Method} & \textbf{Syntactic} & \textbf{Functional} & \textbf{Mean Time} & \textbf{Median} & \textbf{P95} & \textbf{Constraint \%} \\
\hline
NO-CD   & 85.4\% & 46.9\% & 9.48s  & ---   & ---    & 0\%   \\
IG-CD   & 94.9\% & 52.1\% & 8.89s  & 8.55s & 13.51s & 18.1\% \\
LAVE    & \textbf{98.9\%} & \textbf{55.0\%} & 7.60s  & 4.45s & 23.59s & 8.5\%  \\
\hline
\end{tabular}
\end{table*}

Our reproduction results match the originally reported numbers closely. The small discrepancies in median and tail latency (P95, Max) are attributable to non-deterministic I/O scheduling on shared GPU resources and the stochastic nature of LAVE’s sampling-based verification, since the original paper does not fix a random seed for the lookahead sampling procedure.

\paragraph{Evaluation on JSONSchemaBench.}
Because JSON-Bench uses relatively simple schemas, we additionally evaluate LAVE on JSONSchemaBench \citep{geng2025jsonschemabenchrigorousbenchmarkstructured}, a harder benchmark that includes deeply nested, recursive, and cross-referential JSON schemas. According to Table~\ref{tab:main-comparison}, LAVE already establishes state-of-the-art performance over IG-CD on the simpler tasks. On this harder dataset, IG-CD is therefore not included, as our experiments found it to require substantially longer inference time and higher cost while offering no performance advantage over LAVE. Thus, we report only LAVE results in Table~\ref{tab:jsonschemabench}.

\begin{table}[h]
\centering
\caption{LAVE performance on JSONSchemaBench (harder schemas). Skipped instances failed grammar validation before inference; timeouts exceeded the 120s budget.}
\label{tab:jsonschemabench}
\begin{tabular}{lc}
\hline
\textbf{Metric} & \textbf{LAVE} \\
\hline
Skipped (grammar invalid) & 75        \\
Results (attempted)       & 511       \\
Valid                     & 396 / 511 (77.5\%) \\
Timeouts ($>$120s)        & 73        \\
Mean time                 & 27.4s     \\
Median time               & 21.9s     \\
P95 time                  & 73.4s     \\
Max time                  & 113.2s    \\
\hline
\end{tabular}
\end{table}

The degradation from 98.9\% on JSON-Bench to 77.5\% valid on JSONSchemaBench, combined with the sharp increase in latency (mean 7.60s $\to$ 27.4s), reveals a substantial robustness gap that motivates our proposed approach.



\section{Error Analysis}

\paragraph{Failure patterns.}
We systematically examine LAVE’s failures on JSONSchemaBench across three categories:

\begin{enumerate}
    \item \textbf{Grammar-invalid skips (75 instances).} Before inference even begins, 75 instances are skipped because the schema itself cannot be parsed by LAVE’s incremental grammar checker. These are schemas with advanced features such as \texttt{\$ref} cycles, \texttt{anyOf}/\texttt{oneOf} with non-trivial unions, or deeply nested \texttt{allOf} compositions. LAVE’s parser relies on a single-path incremental FSM and cannot handle the branching or backtracking required for these constructs.

    \item \textbf{Timeouts (73 instances).} Instances that exceed the 120-second budget represent cases where LAVE’s lookahead sampling enters a long retry loop. This occurs when a grammar state requires a specific sequence of tokens (e.g., a long mandatory key name or a fixed enum value) that the diffusion model rarely proposes in the sampled completions. Each failed sample triggers a full re-generation of the remaining masked span, and with complex schemas the valid token space at a given position can be vanishingly small. The P95 latency of 73.4s confirms that a significant tail of instances is pathologically slow.

    \item \textbf{Invalid outputs (115 / 511 instances).} Among attempted instances that did not time out, 22.5\% produce invalid JSON. Inspection reveals two sub-patterns: (a) \emph{suffix incompatibility}, where an accepted token at position $i$ is locally valid but makes the suffix unreachable under grammar constraints. For example, committing to a string value where a number is later required by a schema dependency; and (b) \emph{stochastic false acceptance}, where LAVE accepts a token because at least one of the $N$ lookahead samples happens to parse, but the accepted token’s probability mass is concentrated on a non-parsing continuation in practice.
\end{enumerate}

\paragraph{Root causes.}
The core issue is that LAVE performs \emph{independent, position-wise} verification. Each token decision is made by sampling the remaining masked span in isolation, which means: (i) the joint distribution over simultaneously unmasked positions is never considered; and (ii) the parser state carried into the lookahead samples does not account for tokens that will be fixed in later diffusion steps. Consequently, tokens that are valid in a local sense may still produce sequences that are globally invalid, and the stochastic retry mechanism offers only a probabilistic guarantee of correctness rather than a structural one.


\section{Reflection}

\paragraph{What we learned.}
The baseline reproduction confirms that grammar-constrained decoding is essential for dLLMs: NO-CD achieves only 85.4\% syntactic validity on the simpler JSON-Bench, and LAVE’s validity drops to 77.5\% on JSONSchemaBench. More importantly, the error analysis reveals that LAVE’s stochastic, position-wise verification is the fundamental bottleneck, not the diffusion backbone itself.

\paragraph{Capabilities needed.}
A stronger model would require: (a) \emph{joint} multi-position grammar checking that evaluates token combinations rather than individual tokens; (b) a \emph{deterministic} acceptance criterion that does not rely on lucky samples to confirm a token is safe; and (c) a parser backend that supports recursive/branching schemas, not just linear FSM paths.

\paragraph{Strengths of the baseline.}
LAVE’s key strength is its low constraint overhead on simple schemas (8.5\% of tokens affected, mean 7.60s on JSON-Bench), which our beam-search approach may not match due to the cost of maintaining multiple hypotheses per diffusion step. This is a genuine trade-off that GGBS must address.

\paragraph{Refined proposed approach.}
The error analysis directly informs our GGBS design: (i) beam expansion over the full masked span addresses suffix incompatibility by jointly evaluating token assignments; (ii) grammar-guided pruning during expansion (rather than post-hoc verification) eliminates the stochastic false-acceptance problem; and (iii) suffix feasibility checks at the beam boundary handle the right-context incompatibility failures. We additionally plan to incorporate a CSR-compiled automaton backend to support the \texttt{\$ref} and \texttt{anyOf} constructs that caused the 75 grammar-invalid skips.

\paragraph{Fallback directions.}
If GGBS proves too slow for the JSONSchemaBench tail (where schemas require large beam widths to find valid paths), our error analysis suggests a hybrid strategy: use LAVE on simple schemas (where it is fast and reliable) and fall back to GGBS only when the schema complexity exceeds a threshold detectable from the schema structure. A secondary direction is to preprocess complex schemas into a normalized form that LAVE’s FSM parser can handle, bypassing the grammar-invalid skip problem without requiring full beam search.
